\documentclass[a4paper, 11pt]{article}
\usepackage[margin=2cm]{geometry}
\usepackage{graphicx} % Required for inserting images
\usepackage[brazilian]{babel}
\usepackage{tikz}
\usepackage{pgfplots}
\usepackage{xfp}
\usepackage{rmnc}
\usepackage[version=4]{mhchem}
\usepackage{chemfig}
\usepackage{xcolor}
\pgfplotsset{compat=1.18}
\newcommand{\cor}{blue}

\title{Pacote rmnc}
\author{Daniel Sedano Ramilo$^1$}
\date{Abril de 2026}

\setlength{\parindent}{0pt}

\begin{document}
	\maketitle
	{\noindent \small $^1$ Graduação em Química, UFRJ, Instituto de Química, Rio de Janeiro – RJ, \textit{danielsr@gradu.iq.ufrj.br}}
	
	\section{Introdução}
	
	\section{Uso}
	Pode-se usar o pacote informando no preâmbulo do seu documento:
	\begin{verbatim}
		\usepackage{rmnc}
	\end{verbatim}
	
	Para o uso pleno do pacote é necessário compilar o arquivo com a versão \LaTeXe \space e será necessário a declaração dos seguintes pacotes no preâmbulo do arquivo \TeX:
	\begin{verbatim}
		\usepackage{tikz}
		\usepackage{pgfplots} %Geralmente pede-se \pgfplotsset{compat=1.18}
		\usepackage{xfp}
	\end{verbatim}
	
	O Espectro pode ser desenhado iniciando o ambiente \verb|espectroc|. O ambiente por si só gera um retângulo aonde poderá ser adicionados os sinais
	\begin{verbatim}
		\begin{espectroc}[<opções>]
			...
		\end{espectroc}
	\end{verbatim}
	
	\begin{espectroc}
	\end{espectroc}
	
	Os parâmetros indicados em ``{\it opções}'' muda o estilo do espectro base. É importante lembrar que os parâmetros altera o espectro local e não serve para os demais espectros que você desenhar no seu documento. Os parâmetros, opcionais, que podem ser informados são:
	
	\begin{table}[h]
		\begin{tabular}{r c l}
			$<$opção$>$ & Padrão & Descrição\\
			\textit{ppmini}	& 0		& Valor inicial do eixo x.\\
			\textit{ppmfin}	& 200	& Valor final do eixo x.\\
			\textit{escala}	& 50	& Valor do múltiplo que aparecerá no eixo (ex.: de 0; 50; 100; 150; 200)\\
			\textit{dominio}	& 0:200	& Formalização para a linha base, coloque os valores de ppmini:ppmfin, caso altere\\
			\textit{largura} & 16cm	& Largura do espectro, pode-se usar \verb|\linewidth| para usar a largura da página\\
			\textit{altura}	& 5cm	& Altura do espectro
		\end{tabular}
	\end{table}
	\newpage
	
	O código:
	\begin{verbatim}
		\begin{espectroc}[ppmini=50, escala=10, largura=\linewidth]
			...
		\end{espectroc}
	\end{verbatim}
	gera o espectro base:
	
	\begin{espectroc}[ppmini=50,
					  escala=10, 
					  largura=\linewidth]
		...
	\end{espectroc}
	
	Como os espectros de RMN de carbono-13 desacoplados de hidrogênio não aprensentam multiplicidade dos sinais, para colocar um sinal no espectro basta indicar dentro do ambiente \verb|espectroc| o código:
	\begin{verbatim}
		\sinal[int]{ppm}
	\end{verbatim}
	Onde o parâmetro \verb|int| é um valor {\bf inteiro} entre 0 e 100 que diz a intensidade do sinal por padrão o valor é 100. O valor \verb|ppm| indica o valor (pode-se usar casas decimais substituindo a vírgula por ponto Ex.: 20.5) do eixo $x$ onde o sinal será apresentado. Além disso, por algum motivo que eu, o desenvolvedor, não saberia explicar, não é necessário a presença do ponto e vírgula após o comando \verb|\sinal|.
	
	\begin{verbatim}
		\begin{espectroc}[escala=25]
			\sinal[int=50]{20};
			\sinal[int=80]{100}
			\sinal[int=30]{105}
		\end{espectroc}
	\end{verbatim}
	
	\begin{espectroc}[escala=25]
		\sinal[int=50]{20};
		\sinal[int=80]{100}
		\sinal[int=30]{105}
	\end{espectroc}
	
	\newpage
	\section{Exemplos}
	\subsection{Exemplo 1}
	\begin{verbatim}
		\noindent
		\begin{minipage}[t]{0.2\linewidth}
			\vspace{1px}
			Acetato de Etila\\
			\ce{C3H6O2}\par
			\vspace{12pt}
			\chemfig[atom sep=23px]{-[:30](=[2]O)-[:-30]O-[:30]}\par
			
		\end{minipage}
		\begin{minipage}[t]{0.8\linewidth}
			\vspace{1px}
			\begin{center}
				\begin{espectroc}[largura=\linewidth]
					\sinal[int=95]{14}
					\sinal[int=87]{21}
					\sinal{60}
					\sinal[int=48]{171}
					\sinal[int=10]{77}
					\sinal[int=9]{78}
				\end{espectroc}
			\end{center}
		\end{minipage}
	\end{verbatim}
	\par
	
	\noindent
	\begin{minipage}[t]{0.2\linewidth}
		\vspace{1px}
		Acetato de Etila\\
		\ce{C3H6O2}\par
		\vspace{12pt}
		\chemfig[atom sep=23px]{-[:30](=[2]O)-[:-30]O-[:30]}\par
		
	\end{minipage}
	\begin{minipage}[t]{0.8\linewidth}
		\vspace{1px}
		\begin{center}
			\begin{espectroc}[largura=\linewidth]
				\sinal[int=95]{14}
				\sinal[int=87]{21}
				\sinal{60}
				\sinal[int=48]{171}
				\sinal[int=10]{77}
				\sinal[int=9]{78}
			\end{espectroc}
		\end{center}
	\end{minipage}
	
	\newpage
	\subsection{Exemplo 2}
	\begin{verbatim}
		Para a molécula de pent-3-ona temos três sinais sendo o $\alpha$ o sinal do carbono
		da carbonila, o $\beta$ o sinal do carbono ligado ao carbono da carbonila
		e o $\gamma$ o sinal do carbono da extremidade.
		\begin{center}
			\begin{espectroc}[largura=\linewidth,
				dominio=0:250,
				escala=25,
				ppmini=0,
				ppmfin=250]
				\sinal[int=90]{8}
				\sinal{35.5}
				\sinal[int=40]{211}
				
				\node[anchor=south] at (axis cs:211,40) {$\alpha$};
				\node[anchor=east] at (axis cs:35,90) {$\beta$};
				\node[anchor=east] at (axis cs:8,80) {$\gamma$};
			\end{espectroc}
		\end{center}
	\end{verbatim}
	
	
	Para a molécula de pent-3-ona temos três sinais sendo o $\alpha$ o sinal do carbono da carbonila, o $\beta$ o sinal do carbono ligado ao carbono da carbonila e o $\gamma$ o sinal do carbono da extremidade.
	\begin{center}
		\begin{espectroc}[largura=\linewidth,
			dominio=0:250,
			escala=25,
			ppmini=0,
			ppmfin=250]
			\sinal[int=90]{8}
			\sinal{35.5}
			\sinal[int=40]{211}
			
			\node[anchor=south] at (axis cs:211,40) {$\alpha$};
			\node[anchor=east] at (axis cs:35,90) {$\beta$};
			\node[anchor=east] at (axis cs:8,80) {$\gamma$};
		\end{espectroc}
	\end{center}
	
	
\end{document}